\documentclass[11pt,letterpaper]{article}

\usepackage[T1]{fontenc}
\usepackage{lmodern}
\usepackage[margin=1in]{geometry}
\usepackage{amsmath,amssymb}
\usepackage{graphicx}
\usepackage{booktabs}
\usepackage[hidelinks]{hyperref}

% Preserve readable text when copying or searching the pdfLaTeX output.
\input{glyphtounicode}
\pdfgentounicode=1

\hypersetup{pdftitle={Simulation-Based Inference: STA492 Project Report}}

% Replace each visible TODO prompt with your own writing.
% The section structure is a starting point; adapt it to your research question.
\newcommand{\placeholder}[1]{\textit{[TODO: #1]}}

\title{Simulation-Based Inference\\[0.3em]
  \large STA492 Project Report}
\author{\placeholder{Group member names}}
\date{\placeholder{Submission date}}

\begin{document}

\maketitle

\begin{abstract}
\placeholder{Summarize the research problem, the assigned papers, the directions
explored, and your main judgments about the most promising opportunities. Write
this after developing the main argument.}
\end{abstract}

\section{Introduction}
\label{sec:introduction}

\placeholder{Introduce the scientific and statistical problem. Explain why
simulation-based inference is relevant, state the focus of your reading, and
describe the questions that guide this report.}

\section{Statistical Background}
\label{sec:background}

\placeholder{Define the statistical setting, notation, and assumptions needed to
understand the papers. Explain the central inferential challenge and the methods
most relevant to your research question.}

% Example equation: uncomment and adapt when useful.
% \begin{equation}
%   p(\theta \mid x) \propto p(x \mid \theta)\,p(\theta).
%   \label{eq:posterior}
% \end{equation}
% Refer to it with \eqref{eq:posterior}.

\section{Critical Review of the Assigned Papers}
\label{sec:papers}

% Add or remove paper subsections to match the assigned reading.
\subsection{Assigned Paper 1}

\placeholder{Identify and cite the paper. Explain its research question, main
contribution, statistical ideas, assumptions, evidence, and limitations. Distinguish
the authors' claims from your own assessment.}

\subsection{Assigned Paper 2}

\placeholder{Identify and cite the paper. Explain its research question, main
contribution, statistical ideas, assumptions, evidence, and limitations. Distinguish
the authors' claims from your own assessment.}

\subsection{Synthesis and Open Questions}

\placeholder{Connect the papers. Discuss where their ideas agree or differ,
which limitations matter most, and which unresolved questions motivate further
research.}

\section{Exploration of Research Directions}
\label{sec:directions}

\placeholder{Describe how the assigned papers informed the five AI-proposed
directions below. Evaluate scientific value, statistical soundness, novelty,
feasibility, and whether each proposal is worth pursuing. Explain how the group
checked, challenged, or revised the suggestions.}

% Give each direction a descriptive title when the proposal is available.
% For each, cover the question, rationale, statistical assumptions, relevant
% literature, practical constraints, and your group's reasoned judgment.
% The five directions do not need equal space in the completed report.
\subsection{Direction 1}

\placeholder{Describe and evaluate the first proposal, then justify your judgment.}

\subsection{Direction 2}

\placeholder{Describe and evaluate the second proposal, then justify your judgment.}

\subsection{Direction 3}

\placeholder{Describe and evaluate the third proposal, then justify your judgment.}

\subsection{Direction 4}

\placeholder{Describe and evaluate the fourth proposal, then justify your judgment.}

\subsection{Direction 5}

\placeholder{Describe and evaluate the fifth proposal, then justify your judgment.}

% Optional: uncomment if the group develops additional directions.
% \subsection{Directions Developed by the Group}
% \placeholder{Present additional ideas and explain how they arose and how the
% group evaluated them.}

\section{Most Promising Directions}
\label{sec:promising}

\placeholder{Develop the direction or directions that merit closer attention.
Explain why they are more promising than the alternatives, refine the research
questions, and identify useful next steps. Distinguish proposed work from any
exploratory work actually completed.}

% If you include computational work, explain the setup, inputs, software,
% random seeds, and commands needed to reproduce the results. Place code and
% lengthy technical details in supplementary material when appropriate.
% A completed implementation or results section is not a course requirement.

% Example table: uncomment, adapt, and replace all placeholders with evidence.
% \begin{table}[tbp]
%   \centering
%   \caption{Comparison of selected research directions.}
%   \label{tab:directions}
%   \begin{tabular}{lll}
%     \toprule
%     Direction & Main opportunity & Main limitation \\
%     \midrule
%     TODO & TODO & TODO \\
%     \bottomrule
%   \end{tabular}
% \end{table}
% Refer to it with Table~\ref{tab:directions}.

% Example figure: add the image file before uncommenting.
% \begin{figure}[tbp]
%   \centering
%   \includegraphics[width=0.8\linewidth]{figures/example.pdf}
%   \caption{Describe the figure and its relevance to the argument.}
%   \label{fig:example}
% \end{figure}
% Refer to it with Figure~\ref{fig:example}.

\section{Discussion}
\label{sec:discussion}

\placeholder{Discuss the broader implications of your reading and evaluation.
Address uncertainty, limitations of your assessment, and the human judgment
needed to verify AI-generated suggestions.}

\section{Conclusion}
\label{sec:conclusion}

\placeholder{State the main lessons from the papers and the exploration of
research directions. Summarize the strongest opportunities and the reasons for
your final judgment.}

% Replace the unnumbered placeholder below with real \bibitem entries.
% Cite a source with \cite{paper-key} once its entry exists.
\begin{thebibliography}{9}
\item[] \placeholder{Add complete references for the assigned papers and all
other cited sources.}

% \bibitem{paper-key}
% Author(s).
% \emph{Paper title}.
% Journal or conference, year.
% \url{https://doi.org/REPLACE-WITH-DOI}
\end{thebibliography}

% Optional supplementary appendix: uncomment when supporting material is ready.
% \clearpage
% \appendix
% \section{Supplementary Material}
% \label{app:supplement}
% Add derivations, detailed evaluations, or reproducibility instructions that
% support the main argument. Cite this appendix in the relevant main-text section.

\end{document}
